% ARCHIVED FROM v82 EXACT SOURCE EXCERPT: 02_fixed-h-mcfg/japanese/main_JP.tex
% Source reference: archive/mcfg-v82-pretrim-full-20261018 (commit bcc0caf0eade66f552127171c0de24d2b92950b0)
% Restoration: replace the corresponding trimmed location with this excerpt (without these header lines).

\section{Yoshinaka の一般枠組みとの関係}
\label{sec:general-framework}

規則ランク有界の学習器は，Yoshinaka の代入可能な
signature に関する枠組み~\cite{Yoshinaka2015General} の
型付きの特別な場合である．
この対応を明示するのは，先行研究の質的な学習結果と，
第~\ref{sec:char-data}節の総文字数上界とを区別するためである．

有限アルファベット $\Sigma$，$f,r\ge1$ および
$h:\Sigma^*\to M$ を固定し，
$M^+:=h(\Sigma^+)$ とおく．
signature $\mathfrak S_{f,r,h}$ は，
開始ソート $O_*=\Sigma^*$ と，
タグ付きソート
\[
 T_{d,\boldsymbol\mu}
 =\{\tuple w\in(\Sigma^+)^d:
 h^{(d)}(\tuple w)=\boldsymbol\mu\},
 \qquad
 1\le d\le f,\quad\boldsymbol\mu\in(M^+)^d
\]
をもつ．
その関数記号は，子の規則ランクが高々 $r$ である
型の整合したすべての良形 MCFG テンプレート
（ランク0を含む），
昇格演算 $T_{1,\mu}\to O_*$，
および $O_*$ に属する零項の $\eps$ 定数からなる．
$O_*$ を入力として受け取る関数はない．
型のタグを忘れると，これは
~\cite[Section~8]{Yoshinaka2015General}
の意味で，非空タプル上の良形
$\lambda$-自由 MCFG 演算の signature を
有限ソート精密化したものである．
この型なし signature は
~\cite[Example~3]{Yoshinaka2015General}
の例と文字どおり同一ではなく，
\emph{非空成分に制限した変種}である．
ファンアウト1における接頭辞--接尾辞型付けの場合は，
~\cite[Example~4]{Yoshinaka2015General}
に対応する．
この signature 上のある文法で実現可能な
部分構造と文脈を，それぞれ
$\mathcal S_O,\mathcal C_O$ と書く．

\begin{lemma}[Signature の文脈は許容文文脈である]
\label{lem:signature-contexts}
開始ソート以外の各ソートについて
\[
 \mathcal S_{T_{d,\boldsymbol\mu}}
 =T_{d,\boldsymbol\mu},
 \qquad
 \mathcal C_{T_{d,\boldsymbol\mu}}
 =\{\tuple x\mapsto E[\tuple x]:E\in\mathcal E_d\}
\]
が成り立つ．
後者の等式は $T_{d,\boldsymbol\mu}$ 上の
\emph{関数}の集合の等式である．
また $\mathcal C_{O_*}$ に含まれるのは
恒等関数だけである．
\end{lemma}

\begin{proof}
開始ソート以外の各ソートの要素は，
ランク0の良形テンプレート定数として表現できる．
$O_*$ を入力に取る演算は存在しない．
開始ソート以外のソートに対する文脈について，
穴から開始記号へ至る経路に沿って帰納法を行う．
良形テンプレートと許容される親文脈を合成しても，
区別した子の変数の出現順序と
各変数の一回限りの出現は保存される．
同じ出力成分内の連続する穴は，
終端記号または非空の兄弟成分で分離される
（非融合性）．
別々の出力成分にある穴は，
非空の親文脈の区切り語によって分離される．
したがって補題~\ref{lem:filling-identity} を
繰り返し適用すれば，
得られるすべての文文脈は $\mathcal E_d$ に属する．

逆に，任意の
$E=v_0\blank_1v_1\cdots\blank_dv_d
\in\mathcal E_d$ に対し，
単項の良形テンプレート
$(v_0\xi_1^1v_1\cdots\xi_1^dv_d)$ を，
ランク0の子と開始ソートへの昇格演算とともに用いれば，
関数 $\tuple x\mapsto E[\tuple x]$ が実現する．
内部区切り語 $v_i$ は良形性により非空であり，
出力の型は $h$ によって決まる．
\end{proof}

\begin{proposition}[Yoshinaka の2015年の枠組みによる規則ランク有界クラスの包含]
\label{prop:yoshinaka2015-embedding}
signature $\mathfrak S_{f,r,h}$ は
~\cite{Yoshinaka2015General} の
Assumptions~1--3 を満たし，
\(\mathsf{Subst}(\mathfrak S_{f,r,h}) =\Cmcf{f}{h}\cap\mathsf{MCFL}(f,r)\)
となる．
したがって，規則ランク有界のクラスには
Yoshinaka の同定定理
~\cite[Theorem~1]{Yoshinaka2015General}
を適用できる．
$f,r,h$ を固定すれば，
蓄積標本に関する更新時間は多項式であり，
刈り込まれた signature 文法表示について，
特徴集合の要素数は
$|N^\sharp|+|P^\sharp|$ 以下である
~\cite[Lemma~3]{Yoshinaka2015General}．
ただし，これは特徴語の\emph{総文字数}を
抑えるものではない．
\end{proposition}

\begin{proof}
signature 文法の型タグを忘れれば，
ファンアウトが高々 $f$，
規則ランクが高々 $r$ の MCFG になる．
逆に，ランク保存の良形標準化
（命題~\ref{prop:good-normalization}），
刈り込んだ生成語型精密化
（定義~\ref{def:output-refinement}），
開始ソートへの昇格演算と独立の $\eps$ 定数によって，
$\mathsf{MCFL}(f,r)$ のすべての言語を
この signature 上で表現できる．
開始ソートでは代入可能性の条件は自明である．
それ以外のソートでは，
補題~\ref{lem:signature-contexts} により，
signature における四例の含意は，
各アリティと成分別型について
定義~\ref{def:tuple-substitutable} と一致する．
刈り込んだ表示には不要な非終端記号や
関数が含まれないので，
~\cite[Section~6]{Yoshinaka2015General}
の要求も満たす．
これで表示した言語クラスの等式が従う．

実効性に関する仮定については，
ソートや文脈への所属判定と
テンプレートの評価が，
上の補題と有限型付けによって多項式時間で行える．
規則ランク有界の MCFG の一様な所属判定も
多項式時間である
~\cite[Proposition~4]{Yoshinaka2011}．
すべての演算のアリティは高々 $r$ である．
長さ $n$ の文について，
開始語，非空の内部区切り語をはさみ
互いに重ならない高々 $f$ 個の非空区間による
タグ付きタプル，および対応する文脈関数を列挙する．
各親タプル出現について，
子の成分に対応する順序付き区間
（高々 $fr$ 個）と残りの終端部分で決まる
ランク高々 $r$ の良形テンプレートを列挙する．
ここでは実際の\emph{関数記号}を記録する．
必要に応じて昇格演算と $\eps$ 定数も加える．
補題~\ref{lem:composition-witness-materialization} により，
実現可能な関数のインスタンスはすべて列挙に現れる．
必要な切断位置は高々 $2f(r+1)$ 個なので，
$f,r,h$ を固定すれば列挙は $n$ の多項式時間である．
これらはまさに Assumptions~1--3 に対応する．
したがって Yoshinaka の Theorem~1 と Lemma~3 を適用すると，
主張した学習に関する帰結が得られる．
\end{proof}

\begin{lemma}[部分構造添字付き再構成と型制約付き再構成]
\label{lem:substp}
上の signature と列挙を用いると，
任意の有限正例標本 $K\subseteq\Sigma^*$ について，
Yoshinaka の $\mathrm{SUBST}_{\mathrm P}(K)$ と
$\mathcal B^{\le r}_{f,h}(K)$ は
同じ言語を生成する．
\end{lemma}

\begin{proof}
まず同じ観測タプルと型付きソートで添字づけられた
開始記号以外の非終端記号を対応させる．
親 $s_0$ とテンプレート値
$z=\rho(s_1,\ldots,s_t)$ をもつ
$\mathrm{SUBST}_{\mathrm P}$ の規則には，
~\cite{Yoshinaka2015General} の記法で
\(c\in\mathrm{CON}_O(K),\qquad c\odot s_0,\ c\odot z\in K\)
を満たす共通文脈 $c$ が必要である．
補題~\ref{lem:signature-contexts} により，
これは許容される文文脈である．
型制約付き Rule~(Link) は
$[s_0]$ から $[z]$ へ遷移する．
さらに Rule~(Comp)，または $t=0$ なら
Rule~(Const) が続く．
その証拠の存在は
補題~\ref{lem:composition-witness-materialization}
で保証される．
逆に，型制約付きの各 Link は，
単項恒等テンプレートに対する
$\mathrm{SUBST}_{\mathrm P}$ のインスタンスであり，
各 Comp と Const は
それ自体が親の出現によって裏づけられたインスタンスである．
したがって導出に関する帰納法により，
開始記号以外の対応する非終端記号が
生成する言語は一致する．

開始ソートにおける文脈は恒等関数だけであり，
開始ソートを引数に取る演算は存在しない．
したがって観測された各開始記号は，
昇格演算と直接の $\eps$ 定数によって，
同じ言語
\(\bigcup_{w\in K\cap\Sigma^+}L([(w)]) \ \cup\ (\{\eps\}\cap K)\)
を生成する．
これは型制約付きの Rule~(Start) と正確に一致する．
$K=\emptyset$ の場合は両構成とも空言語を生成する．
\end{proof}

\paragraph{特徴証拠の対応．}
Yoshinaka の証明では文脈添字付きの
$\mathrm{SUBST}_{\mathrm D}$ を使用するが，
同論文の第7節では追加の仮定の下で
$\mathrm{SUBST}_{\mathrm P}$ についても論じている
~\cite{Yoshinaka2015General}．
本稿の標本ごとの言語の一致は，
補題~\ref{lem:substp} によって直接証明した．
開始ソート以外の証拠を
$c_X=\chi(X)$ と $s_X=\omega(X)$
として選ぶと，
$\Wit(\widetilde G_0)$ と正確に対応する．
開始ソートへの昇格に関する証拠は
選ばれた語の中に既に含まれ，
$\eps$ が必要な場合には
その直接の証拠を加える．
ここで同一視しているのは，
開始ソートの規約をそろえたときの
正例語からなる証拠の\emph{構成パターン}であり，
仮説文法の構文的な同一性を主張するものではない．


